\documentclass{article}

\usepackage[T1]{fontenc}
\usepackage[spanish, es-nodecimaldot, es-noquoting, es-tabla]{babel}
\usepackage[letterpaper,top=2cm,bottom=2cm,left=3cm,right=3cm,marginparwidth=1.75cm]{geometry}

\usepackage{amsmath}
\usepackage{xcolor}
\usepackage{listings}
\usepackage{booktabs}
\usepackage[colorlinks=true, allcolors=blue]{hyperref}

\setlength{\emergencystretch}{3em}

% El código resuelto, copiado de Intro2AI/solutions/practice_bayesian_learning_v4
\lstset{
    language=Python,
    basicstyle=\ttfamily\small,
    keywordstyle=\color{blue!70!black},
    commentstyle=\color{green!40!black}\itshape,
    stringstyle=\color{red!60!black},
    backgroundcolor=\color{gray!6},
    frame=single,
    rulecolor=\color{gray!40},
    framesep=4pt,
    xleftmargin=4pt,
    xrightmargin=4pt,
    columns=fullflexible,
    keepspaces=true,
    showstringspaces=false,
    upquote=true,
    aboveskip=6pt,
    belowskip=6pt,
    literate={á}{{\'a}}1 {é}{{\'e}}1 {í}{{\'i}}1 {ó}{{\'o}}1 {ú}{{\'u}}1 {ñ}{{\~n}}1
             {Á}{{\'A}}1 {É}{{\'E}}1 {Í}{{\'I}}1 {Ó}{{\'O}}1 {Ú}{{\'U}}1,
}

\title{Clasificadores Naïve Bayes:\\ las flores iris y un filtro de spam\\[6pt] \large Soluciones (solo para el profesor)}
\author{Arturo Gil y Juan José Cabrera}
\date{}

\begin{document}
\maketitle

\noindent\textit{Este documento no se entrega a los alumnos. Para cada ejercicio, recoge el código resuelto y las respuestas a las preguntas del guion. Todas las cifras se han comprobado ejecutando los programas con los ejercicios resueltos. El código completo está en \texttt{Intro2AI/solutions/practice\_bayesian\_learning\_v4}.}

\section*{Ejercicio 1. Naïve Bayes gaussiano: las flores iris}
En \texttt{naive\_bayes/gaussiano.py}, \texttt{entrenar\_naive\_bayes} (a) y las otras dos funciones (b):
\begin{lstlisting}
def entrenar_naive_bayes(X_train, y_train, num_clases):
    medias, desviaciones, priors = [], [], []
    for c in range(num_clases):
        X_c = X_train[y_train == c]             # solo los ejemplos de la clase c
        medias.append(np.mean(X_c, axis=0))
        desviaciones.append(np.std(X_c, axis=0))
        priors.append(len(X_c) / len(y_train))
    return medias, desviaciones, priors

def log_densidad_gaussiana(x, media, desviacion):
    return (-0.5 * np.log(2 * np.pi * desviacion ** 2)
            - (x - media) ** 2 / (2 * desviacion ** 2))

def calcular_puntuaciones(x, medias, desviaciones, priors, num_clases):
    puntuaciones = {}
    for c in range(num_clases):
        p = np.log(priors[c])
        for i in range(len(x)):                 # atributo a atributo
            p += log_densidad_gaussiana(x[i], medias[c][i], desviaciones[c][i])
        puntuaciones[c] = p
    return puntuaciones
\end{lstlisting}

\begin{enumerate}
    \item[a)] \textbf{El pétalo.} En el largo del pétalo ($x_3$), setosa (por debajo de 2~cm) queda lejísimos de las otras dos variedades, con medias de 4,23 y 5,53~cm, y las campanas de versicolor y virginica solo se solapan un poco, hacia los 5~cm. El ancho del pétalo ($x_4$) separa casi igual de bien. Los dos atributos del sépalo se solapan mucho, sobre todo el ancho ($x_2$), con las tres medias entre 2,77 y 3,44~cm: por sí solos apenas servirían. Los errores se esperan entre \textbf{versicolor y virginica}, cuyas campanas se solapan en los cuatro atributos; setosa no se confunde con nada. Todo lo que el clasificador sabe del problema son 3 probabilidades a priori y 24 números: 12 medias y 12 desviaciones típicas.

    Si las desviaciones típicas no coinciden: con \texttt{np.var} en lugar de \texttt{np.std} salen las varianzas, 0,0276, 0,2074 y 0,2674, y con la desviación típica muestral (\texttt{ddof=1}, que divide entre $n - 1$), 0,1683, 0,4616 y 0,5241. \texttt{GaussianNB} usa la que divide entre $n$.

    \item[b)] \textbf{La cuenta a mano.} Con $\mu = 1{,}4921$ y $\sigma = 0{,}1660$, $\sigma^2 = 0{,}02756$ y $2\pi\sigma^2 = 0{,}17314$:
    \[
    \log P(x_3 = 1{,}4 \mid \textit{setosa}) = -\tfrac{1}{2}\log 0{,}17314 - \frac{(1{,}4 - 1{,}4921)^2}{2 \cdot 0{,}02756} = 0{,}8768 - 0{,}1539 = 0{,}7229 \approx 0{,}723
    \]
    \textbf{Un logaritmo positivo.} $P(x_i \mid c)$ es una \textbf{densidad}, no una probabilidad, y una densidad puede valer más que 1: lo que no puede pasar de 1 es el área bajo la campana, la probabilidad de un intervalo. Cuanto más estrecha es la campana, más alta: la del largo del pétalo de setosa ($\sigma = 0{,}166$) llega hasta $1/\sqrt{2\pi\sigma^2} = 2{,}40$, y la del ancho ($\sigma = 0{,}114$), hasta 3,50. Por eso sus logaritmos (0,723 y 1,145) son positivos, y también la suma, 0,939. Para decidir no importa: las tres clases se comparan con densidades.

    En la tabla se ve también qué decide: con el pétalo, \textit{versicolor} y \textit{virginica} suman $-19{,}5$ y $-14{,}8$, y $-32{,}2$ y $-21{,}3$; con el sépalo, solo entre $-1{,}6$ y $-3{,}4$. La flor tiene un pétalo de setosa.

    \item[c)] \textbf{Versicolor y virginica}, una flor en cada sentido, como se esperaba en a). La virginica $(6{,}0;\ 2{,}2;\ 5{,}0;\ 1{,}5)$, de pétalo corto y estrecho para su variedad, sale versicolor ($P = 0{,}96$), y la versicolor $(6{,}9;\ 3{,}1;\ 4{,}9;\ 1{,}5)$ sale virginica, por poco ($P = 0{,}55$). Setosa no se confunde nunca.

    \textbf{Flores nuevas.} El acierto con las flores de entrenamiento (96,5\,\%) mide lo bien que el modelo se ajusta a lo que ya ha visto, no lo bien que generaliza: es optimista, como lo sería un examen con las mismas preguntas que los ejercicios de clase. Lo que interesa es acertar con flores nuevas, y eso solo lo mide un conjunto que no se ha usado para entrenar. Con solo 36 flores de prueba, cada error mueve el acierto casi 3 puntos: es una medida poco precisa, y con otra partición podría salir distinta.

    \texttt{GaussianNB} suma a las varianzas una cantidad minúscula, \texttt{epsilon} (unos $3 \cdot 10^{-9}$ aquí), para que ninguna sea 0. Por eso sus probabilidades difieren un poco de las nuestras, menos de $3 \cdot 10^{-8}$, pero las predicciones son las mismas.
\end{enumerate}

\section*{Ejercicio 2. La bolsa de palabras}
En \texttt{naive\_bayes/bolsa\_de\_palabras.py}:
\begin{lstlisting}
def construir_vocabulario(mensajes):
    vocabulario = {}
    palabras = set()
    for mensaje in mensajes:
        palabras.update(tokenizar(mensaje))
    for i, palabra in enumerate(sorted(palabras)):  # el orden decide los índices
        vocabulario[palabra] = i
    return vocabulario

def vectorizar(mensajes, vocabulario):
    X = np.zeros((len(mensajes), len(vocabulario)), dtype=np.int32)
    for k, mensaje in enumerate(mensajes):
        for palabra in tokenizar(mensaje):
            if palabra in vocabulario:          # si no está en el vocabulario, se ignora
                X[k, vocabulario[palabra]] += 1
    return X
\end{lstlisting}

\begin{enumerate}
    \item[a)] El vocabulario y la matriz:
    \begin{center}
    \small
    \begin{tabular}{lccccccc}
    \toprule
     & \texttt{call} & \texttt{free} & \texttt{me} & \texttt{now} & \texttt{prize} & \texttt{win} & \texttt{you} \\
    \midrule
    \texttt{'Win a FREE prize now'} & 0 & 1 & 0 & 1 & 1 & 1 & 0 \\
    \texttt{'Call me now'} & 1 & 0 & 1 & 1 & 0 & 0 & 0 \\
    \texttt{'I win, you win'} & 0 & 0 & 0 & 0 & 0 & 2 & 1 \\
    \midrule
    \texttt{'Call now to win a prize'} & 1 & 0 & 0 & 1 & 1 & 1 & 0 \\
    \bottomrule
    \end{tabular}
    \end{center}
    \textbf{Siete palabras.} \texttt{tokenizar} pasa todo a minúsculas (\texttt{Win} y \texttt{win} son la misma palabra) y descarta las palabras de una sola letra (\texttt{a} e \texttt{I}), y cada palabra entra en el vocabulario una sola vez, aunque salga en varios mensajes. \textbf{\texttt{to}} no está en el vocabulario, porque no sale en ninguno de los tres mensajes de entrenamiento, así que se ignora. Si no se ignora, \texttt{vocabulario['to']} da un error (\texttt{KeyError}); y si se usa \texttt{vocabulario.get('to')}, que devuelve \texttt{None}, \texttt{X[k, None] += 1} no da error, pero suma 1 en toda la fila, y el vector sale $(2, 1, 1, 2, 2, 2, 1)$.

    \item[b)] \textbf{Ceros.} Un SMS es corto: tiene, de media, 13,3 palabras distintas, de las 7.498 del vocabulario. El resto de su fila son ceros: solo el 0,18\,\% de la matriz (55.546 de sus 31,3 millones de casillas) no es 0.

    \textbf{Matriz dispersa.} Se guardan solo los valores distintos de 0 y, de cada uno, su columna, más dónde empieza cada fila (es el formato CSR, \textit{compressed sparse row}, el de \texttt{CountVectorizer}): dos vectores de 55.546 números (los valores y sus columnas) y uno de 4.181 (dónde empieza cada fila), en lugar de 31,3 millones de casillas. Ocupa unas 180 veces menos, 0,68~MB frente a 125,4~MB. Con un corpus de correo real, de un millón de mensajes y 50.000 palabras, la matriz entera ocuparía 200~GB con 4 bytes por casilla: no cabría en memoria.
\end{enumerate}

\section*{Ejercicio 3. Naïve Bayes multinomial: un filtro de spam}
En \texttt{naive\_bayes/multinomial.py}, la función del entrenamiento (a) y la de las puntuaciones (b):
\begin{lstlisting}
def entrenar_naive_bayes_multinomial(X_train, y_train, num_clases, alpha=1.0):
    num_docs, num_palabras = X_train.shape
    log_priors = np.zeros(num_clases)
    log_prob_palabras = np.zeros((num_clases, num_palabras))
    for c in range(num_clases):
        X_c = X_train[y_train == c]                 # los mensajes de la clase c
        log_priors[c] = np.log(X_c.shape[0] / num_docs)
        conteo_palabras = X_c.sum(axis=0)           # N(palabra, c) de cada palabra
        total_palabras_c = conteo_palabras.sum()    # N(c)
        log_prob_palabras[c] = np.log((conteo_palabras + alpha) /
                                      (total_palabras_c + alpha * num_palabras))
    return log_priors, log_prob_palabras

def calcular_puntuaciones_multinomial(test_sample, log_priors, log_prob_palabras):
    num_clases = len(log_priors)
    num_palabras = len(test_sample)
    puntuacion = np.zeros(num_clases)
    for c in range(num_clases):
        puntuacion[c] = log_priors[c]
        for i in range(num_palabras):       # palabra a palabra del vocabulario
            puntuacion[c] += test_sample[i] * log_prob_palabras[c][i]
    return puntuacion
\end{lstlisting}
Para entrenar no hace falta recorrer las palabras: \texttt{X\_c.sum(axis=0)} da a la vez el conteo de las 7.498, y las operaciones con ese vector (sumarle $\alpha$, dividirlo, \texttt{np.log}) se hacen palabra a palabra.

\begin{enumerate}
    \item[a)] Con $\alpha = 1$, $N(\texttt{ham}) = 47.209$ y $N(\texttt{spam}) = 13.037$ palabras:
    \[
    P(\textit{prize} \mid \texttt{ham}) = \frac{0 + 1}{47.209 + 7.498} = 1{,}83 \cdot 10^{-5}
    \qquad
    P(\textit{prize} \mid \texttt{spam}) = \frac{77 + 1}{13.037 + 7.498} = 0{,}0038
    \]
    cuyos logaritmos son $-10{,}9097$ y $-5{,}5732$. \textbf{Sin suavizado}, $P(\textit{prize} \mid \texttt{ham}) = 0/47.209 = 0$, y como la verosimilitud es un producto, un mensaje legítimo que dijera \texttt{prize} tendría $P(\texttt{ham} \mid \text{mensaje}) = 0$: sería spam con total seguridad, dijera lo que dijera el resto del mensaje. Una sola palabra anula una clase, y el 88\,\% de las palabras (6.626 de 7.498) no sale nunca en una de las dos. Con el código, aún es peor: $\log 0 = -\infty$, y en el bucle de \texttt{calcular\_puntuaciones\_multinomial} cada palabra que no sale en el mensaje suma $0 \cdot (-\infty)$, que da \texttt{nan}. Todos los mensajes de prueba acaban con \texttt{nan} en las dos clases, y \texttt{np.argmax} se queda con la primera, \texttt{ham}: el filtro dice siempre \texttt{ham}. No haber visto una palabra en 4.180 SMS no significa que sea imposible, y el suavizado le reserva un poco de probabilidad.

    La comprobación con $\alpha = 0{,}5$ detecta un error que con $\alpha = 1$ pasa inadvertido: sumar en el denominador el tamaño del vocabulario sin multiplicarlo por $\alpha$.

    \item[b)] \textbf{El bucle de las palabras.} Para un mensaje, la suma del bucle interior se ejecuta $2 \times 7.498 = 14.996$ veces, una por clase y palabra del vocabulario (casi 21 millones para los 1.394 mensajes de prueba), pero solo añaden algo las palabras del mensaje, 12,4 distintas de media: para las demás, \texttt{test\_sample[i]} es 0. Bastaría con recorrer las palabras del mensaje, que es lo que hacen las librerías gracias a las matrices dispersas del ejercicio 2.

    \textbf{La tasa de acierto engaña} porque las clases están descompensadas: el 86,6\,\% de los SMS de prueba son \texttt{ham} (1.207 de 1.394), así que acertar mucho es fácil, y la tasa de acierto apenas dice nada de la clase pequeña, que es la que interesa. Hay que mirar la matriz de confusión, o la precisión y la exhaustividad: $166/(166 + 8) = 0{,}954$ y $166/(166 + 21) = 0{,}888$. Es el mismo fenómeno que el de una prueba médica de una enfermedad rara: cuando casi todos los casos son de una clase, la probabilidad a priori pesa mucho.

    \textbf{Los falsos positivos} son más graves para el usuario: un mensaje legítimo en la carpeta de spam puede no leerse nunca (una oferta de trabajo, un aviso del banco), mientras que un spam en la bandeja de entrada se borra en un segundo. Por eso los filtros prefieren una precisión alta, aunque se cuele algo más de spam.

    \item[c)] \textbf{Frecuente no es indicativo.} Las 10 más frecuentes son casi las mismas en las dos clases: las palabras más comunes del inglés (\texttt{you}, \texttt{to}, \texttt{the}\dots). Como son igual de probables en las dos, $\log P(w \mid \texttt{spam}) - \log P(w \mid \texttt{ham})$ es casi 0, y apenas mueven la decisión: lo que distingue es la diferencia. Indican spam \texttt{claim}, \texttt{prize}, \texttt{150p} (el precio de un SMS de tarificación especial), \texttt{uk}, \texttt{tone} (de los tonos de llamada) y los números; indican \texttt{ham} \texttt{gt} y \texttt{lt} (restos de los códigos HTML \texttt{\&gt;} y \texttt{\&lt;}), \texttt{come}, \texttt{he}, \texttt{she}, \texttt{lor} y \texttt{da} (coloquialismos del inglés de Singapur, de donde viene buena parte de los SMS legítimos). El filtro aprende las peculiaridades de su corpus: con SMS en español no serviría.

    % \item[d)] Dos formas, comprobadas con \texttt{filtro.py}:
    % \begin{itemize}
    %     \item \textbf{Palabras que no conoce.} En
    %     \begin{flushleft}\small\texttt{C0ngratulati0ns! U h4ve w0n a FR33 pr1ze. C4ll n0w 2 cl4im}\end{flushleft}
    %     no hay ninguna palabra del vocabulario (\texttt{U}, \texttt{a} y \texttt{2} tienen un solo carácter), así que no cuenta ninguna, y queda $P(\texttt{spam} \mid \text{mensaje}) = P(\texttt{spam}) = 0{,}134$: \texttt{ham}. Es lo que hacen los spammers al escribir \textit{V1AGRA}: una palabra que no se vio al entrenar se ignora.
    %     \item \textbf{Palabras de \texttt{ham}.} El mensaje
    %     \begin{flushleft}\small\texttt{Congratulations! You have won a free prize, call now to claim it. Sorry lor, I will come later da}\end{flushleft}
    %     tiene todo el spam (\texttt{claim} $+5{,}54$, \texttt{prize} $+5{,}34$\dots), pero queda en $P(\texttt{spam}) = 0{,}02$, porque las palabras de \texttt{ham} lo compensan (\texttt{come} $-4{,}13$, \texttt{lor} $-3{,}82$, \texttt{da} $-3{,}72$, \texttt{later} $-3{,}72$, \texttt{sorry} $-2{,}44$). Naïve Bayes suma la evidencia de cada palabra por separado: basta con añadir suficientes palabras de la otra clase.
    % \end{itemize}
    % Es un ejemplo sencillo de \textbf{ataque adversario} a un modelo de aprendizaje automático. Los filtros de verdad se defienden reentrenándose con el spam nuevo y mirando más cosas que las palabras: el remitente, los enlaces\dots
\end{enumerate}

\end{document}
